\section{好み切替条件}\label{app:switch}
本文の条件では，A-1 の好きな色は常に赤であった．
どのような経験が他者への当てはめを強めるかの手がかりとして，A-1 の好きな色をエピソードごとに赤か緑から選び，快・不快を「好きな色 $-$ 嫌いな色」とした条件（好み切替条件）を作った．
規則の一歩の割合は経験50\%条件と同じ 0.5 で，A-1 が赤を好むエピソードは経験50\%条件のデータと全く同じである．
したがって，経験50\%条件との違いは，残りのエピソードで好きな色が緑に替わることだけである．
組1〜3 で段階Iを学習し，経験50\%条件と同じ評価を行った．
この条件は事前登録したうえで行ったが，段階IIは関門の結果により探索として扱った．
% 出典：v14 ログ §5・§17（データの確認：赤好きのエピソードは v11_s1_mixed と全配列で一致，規則の番 0.500），v14 まとめ §3。

自分側の見込みは，好きな色ごとに対称に生まれた．
自己位置の追跡，価値と行動・色の結びつき，自分側の積み上がりの関門はすべての組で合格し，積み上がり $g$（最終5点の平均）は赤を好むエピソードで $+0.093$〜$+0.105$，緑を好むエピソードで $+0.083$〜$+0.104$ であった．
一方，他者側の関門は G5a・G5b とも3組中1組でしか合格せず（G5a は組3，G5b は組1），他者側への当てはめは弱かった．
A-2 の本当の快・不快を入れても，予測の中の A-2 は A-2 の好きな色にも A-1 自身の好きな色にも寄らなかった．
仮説の照合の BA は 0.50〜0.57 で，経験50\%条件を上回った組はなかった．
1ビットモジュールの一致率も経験50\%条件を上回らなかった．
% 出典：v14 まとめ §3，v14 ログ §33・§37・§38・§41。

A-1 の好きな色を赤から緑に替えて評価し，照合の選択が A-1 自身の好きな色の側へずれるか（自分の好みの投影）も調べたが，どのモデルでも見られなかった．
なお，他者側の特徴予測の誤差で選ぶ場合は，process-2 の内部状態に process-1 の入力が入らないので，この検査は構成上必ず0になる．
% 出典：v14 まとめ §3（投影の検査），v14 ログ §37。
